\section{The two scalar estimates}\label{sc:scalar}\label{sc:section}

The matrix argument has reduced the entropy inequality to two estimates.
Proposition~\ref{mat:scalar-absorption} controls the part depending on the
first Gaussian chaos; Proposition~\ref{mat:scalar-pair} controls the remaining
quadratic form, one pair of eigenvalues at a time.  We prove them here.
The functions $q,d,k,b,l$ and the six intervals
$O,P_1,R_1,U_1,V_1,T_1$ are those of the transport section, and the constants
$\alpha,s,\beta,c,\gamma,\epsilon,\kappa$ and functions $p,e,h_0$ are those
of the matrix section.  All finite decimals below denote exact rational
numbers.

We first state the one-variable bounds used in both proofs.  Their verification
is given in Appendix~\ref{cert:verification}; in particular, none of the
bounds in this section is inferred from numerical sampling.  For the
quadratic form in Proposition~\ref{mat:scalar-pair}, set
\begin{equation}\label{sc:auxiliary-functions}
 \begin{gathered}
 D_*=\frac{1}{4(1-\epsilon)},\qquad
 \mathsf D=\frac{D_*}{p},\quad \mathsf L=\frac{l}{\alpha},\quad
 \mathsf E=e+\gamma b
       =p+b(.75-b)-\frac{l^2}{\alpha},\\
 K_0=\mathsf D\mathsf E,\qquad J_0=\mathsf D cb,\\
 A_*=s+1+s/8,\quad
 \eta_*=s\beta+(1+s/8)\gamma,\quad
 \nu_*=s\beta^2+(9s/8)\gamma^2.
 \end{gathered}
\end{equation}
The symbols $\mathsf D,\mathsf L,\mathsf E$ denote scalar functions, to
be distinguished from the matrix residuals $D_r,L_r,M_r$ of the preceding
section.

\begin{lemma}[One-variable bounds]\label{sc:bounds}
For every $a\in\R$,
\begin{equation}\label{sc:p-h-bounds}
 \begin{gathered}
 p(a)>0,\qquad |h_0(a)|\le .67\sqrt{p(a)},\\
 p(a)\ge\begin{cases}(7+a^2)^{-1},&a\le0,\\ .25,&a\ge0,\end{cases}
 \qquad p(a)\ge .46\quad(a\ge4).
 \end{gathered}
\end{equation}
The following interval bounds hold:
\begin{equation}\label{sc:interval-table}
\begin{array}{c|cccccc}
 &b&\mathsf L\le&\mathsf D\le&K_0&J_0\le&\ell_1\\\hline
 O &[0,.09]&.60&.295(7+a^2)&[.37,.56]&.72&.13\\
 P_1 &[.083,.15]&1.01&4.3&[.29,.46]&.72&.13\\
 R_1 &[.14,.23]&1.19&3.00&[.29,.371]&.68&.13\\
 U_1 &[.22,.342]&1.19&1.72&[.29,.398]&.61&.123\\
 V_1 &[.337,.45]&.83&.87&[.31,.4]&.50&.060\\
 T_1 &[.445,.5]&.20&.640&[.32,.4]&.48&.022
\end{array}
\end{equation}
Here $\ell_1$ is the positive number in the last column, and
\begin{equation}\label{sc:local-surplus}
 Q-
 \begin{pmatrix}
 A_*b^2-\eta_*b+\nu_*/4+l^2/\alpha&\nu_*/2-\eta_*b\\
 \nu_*/2-\eta_*b&\nu_*
 \end{pmatrix}
 -\mathsf D
 \begin{pmatrix}\mathsf E\\cb\end{pmatrix}
 \begin{pmatrix}\mathsf E&cb\end{pmatrix}
 \succeq\begin{pmatrix}\ell_1&0\\0&.75\end{pmatrix}.
\end{equation}
Here $Q=\left(\begin{smallmatrix}.43&.13\\.13&2.72\end{smallmatrix}\right)$,
as in Proposition~\ref{mat:scalar-pair}.
Finally, writing
\begin{equation}\label{sc:derivative-functions}
 z=\frac{l'}{l},\qquad h_*=1+3b+kz,\qquad \chi=\frac{l}{d},
\end{equation}
we have
\begin{equation}\label{sc:derivative-table}
\begin{array}{c|ccc}
 &z&h_*&\chi\le\\\hline
 a\le-1&[-.08,.70]&[.8,1.77]&.070\\
 a\ge-1&[-.475,0]&[.8,1.77]&.097
\end{array}
\end{equation}
At an endpoint shared by two intervals, either applicable row may be used.
\end{lemma}

The first line of the lemma supplies a lower bound for the mean of $p$ from
only a second moment.  The matrix surplus in~\eqref{sc:local-surplus}
then provides room to absorb the secant errors in the second proposition.
We treat these two uses separately.

\subsection{Absorption using a second moment}\label{sc:absorption-proof}

\begin{proof}[Proof of Proposition~\ref{mat:scalar-absorption}]
By $0<b<1/2$ and~\eqref{sc:p-h-bounds},
\begin{equation}\label{sc:y-cost}
 s\left((b-\beta)^2+
 \frac{\alpha h_0^2}{4(1-\epsilon-\kappa)p}\right)
 \le s\left(.28^2+\frac{.075\cdot .67^2}{4\cdot .36}\right)<.111.
\end{equation}
It remains to bound the terms involving $T-\log S$ while retaining a
negative multiple of the weight $\omega$ in
\eqref{mat:rigidity-weight}.

Fix a unit eigenvector $v$ of $B$ with eigenvalue $\lambda\ge0$.
The corresponding eigenvalues of $S,T,H,\rho$
are
\[
 S=(1/s+\alpha\lambda)^{-1},\quad T=S-1,\quad
 H=\alpha S(\lambda-1)^2,\quad \rho=\epsilon+\kappa S/s.
\]
Let $\nu_v$ be the probability measure obtained by averaging the spectral
measure of $A$ at $v$.  The decomposition of $A$ into its first chaos and
$Y$ gives
\begin{equation}\label{sc:spectral-moments}
 \int a\,d\nu_v(a)=0,\qquad
 \int a^2\,d\nu_v(a)=\lambda+r,
 \qquad r=v^{\mathsf T}\E(Y^2)v\ge0.
\end{equation}
Indeed, orthogonality of the Gaussian chaoses gives
$\E A^2=B+\E Y^2$, including the vanishing of both matrix cross terms.

For $0<u\le.13$, define
\[
 n_u=\frac{u^{3/2}}{2},\qquad
 d_u=\frac{u^3}{16(.46-u)}.
\]
We claim the following quadratic minorant:
\begin{equation}\label{sc:p-minorant}
 p(a)\ge u+n_ua-d_ua^2\qquad(a\in\R).
\end{equation}
For $a\le0$, put $x=-\sqrt u\,a$.  If $x\ge2$, the right side is
nonpositive.  If $0\le x\le2$, it is enough to use
$p(a)\ge u/(x^2+7u)$ and
\[
 1-(1-x/2)(x^2+7u)
 \ge \frac{x(x-1)^2}{2}+.045(2-x)\ge0.
\]
For $0\le a\le4$, the right side of~\eqref{sc:p-minorant} is at most
$.13+2(.13)^{3/2}<.25$.  For $a\ge4$, its maximum over all real $a$ is
$u+n_u^2/(4d_u)=.46$.  Thus~\eqref{sc:p-minorant} follows from
\eqref{sc:p-h-bounds} in all three ranges.

Choose
\[
 u=\min(.13,1.28/\sqrt\lambda),
\]
with $u=.13$ when $\lambda=0$.  Since $\lambda u^2\le1.28^2$, integration of the
minorant against~\eqref{sc:spectral-moments} yields
\begin{equation}\label{sc:p-mean}
 \int p\,d\nu_v\ge u-d_u(\lambda+r)
 \ge .689u-d_ur.
\end{equation}
The second inequality uses
$1-1.28^2/[16(.46-.13)]>.689$.
We also have
\begin{equation}\label{sc:rho-loss}
 H\rho d_u\le .027.
\end{equation}
To check this, note that $\rho\le.64$ and that $H\le\lambda$ for $\lambda\ge1$,
whereas $H\le\alpha s$ for $\lambda\le1$.  In the former case,
\[
 H\rho d_u\le
 \frac{.64\cdot1.28^2\cdot.13}{16(.46-.13)}<.02582;
\]
in the latter, $\alpha s\cdot.64\cdot .13^3/[16(.46-.13)]<.027$.

For completeness, we next verify the elementary logarithmic bound needed
to compare $T-\log S$ with the positive term in~\eqref{sc:p-mean}.
Set
\[
 F_0(S)=\frac{S(S-1-\log S)}{(1-S)^2},\qquad F_0(1)=\tfrac12.
\]
For $0<S\le s$ and $\lambda=(1/S-1/s)/\alpha$, we claim
\begin{equation}\label{sc:log-bound}
 \frac{\alpha F_0(S)}{\epsilon+\kappa S/s}
 \le\min(.089,.88/\sqrt\lambda)\le .688u,
\end{equation}
where the second bound in the minimum is interpreted as $+\infty$ at
$\lambda=0$.  First,
\[
 F_0(S)=\int_0^1 \frac{Sx}{1-x+xS}\,dx
\]
is concave.  Its tangent at $S=.28$ has intercept below $.16$ and slope
below $.52$, so $F_0(S)\le .16+.52S$.  These two strict inequalities
can, for example, be checked from
\[
 -\log(.28)=2\sum_{j=0}^{\infty}\frac{(9/16)^{2j+1}}{2j+1},
\]
using twelve terms and the geometric bound for the remainder.
Coefficient comparison now gives
\[
 .075(.16+.52S)<.089(.15+.45325S),
\]
which proves the first bound in the minimum.

For the second bound we use
\begin{equation}\label{sc:F0-sqrt}
 \frac{F_0(S)}{\sqrt S}\le .478+1.4S\qquad(S>0).
\end{equation}
Here is a direct verification, including the point $S=1$ by continuity.
For $x>0$, let
\[
 J(x)=\frac{(.478+1.4x^2)(1-x^2)^2}{x}
       -(x^2-1-2\log x).
\]
Then
\[
 x^2J'(x)=P(x):=-.478+2x+.444x^2-2x^3-6.966x^4+7x^6.
\]
Descartes' rule of signs gives at most three positive roots of $P$,
counting multiplicity.  The signs at $.24,.26,.6$, together with
$P(1)=0$ and $P'(1)>0$, give exactly three: one in $(.24,.26)$, one in
$(.26,.6)$, and $1$.  Thus the only local minima of $J$ occur at the
first root and at $1$.  We have $J(1)=0$ and $J(.25)>.14$; the latter
follows already from $\log2<.694$.  On $[.24,.26]$,
$|P(.25)|<.01$ and $|P'|<4$, hence $|J'|<2$.
The first minimum is therefore positive.  Also $J(x)\to+\infty$ at
both ends of $(0,\infty)$.  This proves $J\ge0$, which is equivalent to
\eqref{sc:F0-sqrt} away from $x=1$.

Using $\alpha\lambda=1/S-1/s$, inequality~\eqref{sc:F0-sqrt} gives
\[
 \frac{\alpha F_0(S)\sqrt\lambda}{\rho}
 \le \sqrt\alpha\sqrt{1-S/s}\,
       \frac{.478+1.4S}{.15+.45325S}
 \le \frac{\sqrt{.075}\cdot .478}{.15}<.88.
\]
The last ratio decreases in $S$ because
$1.4\cdot.15<.478\cdot.45325$.  Finally,
$.089<.688\cdot.13$ and $.88<.688\cdot1.28$ prove the last inequality
in~\eqref{sc:log-bound}.

For $\lambda\ne1$ we have
\[
 \frac{T-\log S}{H\rho}=\frac{\alpha F_0(S)}{\rho}.
\]
Combining~\eqref{sc:p-mean}--\eqref{sc:log-bound}, and treating $\lambda=1$
by continuity, gives
\[
 T-\log S-H\rho\int p\,d\nu_v
       \le -.001H\rho u+.027r.
\]
The definitions give $H\rho u=\omega(\lambda)$.  Sum over an orthonormal
eigenbasis of the matrix $B$, and add~\eqref{sc:y-cost}.
This summation evaluates $\tau(H\rho p(A))$ without requiring $p(A)$
to commute with $B$.  Likewise, \eqref{sc:y-cost} may be integrated
against the positive semidefinite matrix $Y^2$ without a commutation
assumption.  Since $\sum_v r=\norm{Y}^2$, the result is
\[
 (.027+.111)\norm{Y}^2-.001\sum_i\omega(\lambda_i)
 =\frac{69}{500}\norm{Y}^2-\frac1{1000}\sum_i\omega(\lambda_i),
\]
as asserted.
\end{proof}

\subsection{A quadratic estimate for two endpoints}\label{sc:pair-proof}

\begin{proof}[Proof of Proposition~\ref{mat:scalar-pair}]
Fix two endpoints $a_-\le a_+$ and write
\[
 \delta=\log q(a_+)-\log q(a_-),\qquad m_*=m_0(a_-).
\]
When the endpoints coincide, all secant increments vanish, and
\eqref{sc:local-surplus} immediately proves the assertion.
We henceforth suppose $\delta>0$.

We first isolate the cost of the secant increments at one endpoint $a_j$.
Use coordinates $(G,h)$ with $h=g-G/2$.  With $Z=0$, the left side of
the endpoint quadratic form in Proposition~\ref{mat:scalar-pair} is
exactly the two matrices subtracted from $Q$ in~\eqref{sc:local-surplus}.
For an increment vector $(x,y,w)$, define
\begin{align}
 \mathcal N_j(x,y,w)^2={}&A_*x^2+y^2/\alpha+\mathsf D_jw^2\notag\\
 &+\frac{((A_*b_j-\eta_*/2)x+\mathsf L_jy+K_{0j}w)^2}{\ell_{1j}}
   +\frac{(-\eta_*x+J_{0j}w)^2}{.75}.
 \label{sc:increment-norm}
\end{align}
Completing squares against the surplus
$\ell_{1j}G^2+.75h^2$ in~\eqref{sc:local-surplus} shows that the full
endpoint form is bounded above by
\begin{equation}\label{sc:one-endpoint}
 Q(G,h)+Z^2\mathcal N_j(\Delta k,\Delta b,\psi_j)^2.
\end{equation}
Indeed, the two coefficients of $2GZ$ and $2hZ$ are precisely the
numerators of the two squared terms in~\eqref{sc:increment-norm}.
The first line of that formula is the coefficient of $Z^2$ before
completion of squares.

The next step estimates these increment norms by integrating derivatives
between the endpoints.  Recall that a secant slope is the uniform average
of the derivative over $[a_-,a_+]$.  With $\xi=\log q(a)$, differentiation
therefore gives the vector
\begin{equation}\label{sc:derivative-vector}
 v_j(a)=\chi(a)(1,z(a),W_j(a)),\qquad
 W_j(a)=h_*(a)-2b_j-2\mathsf L_jz(a).
\end{equation}
More explicitly, $(\Delta k,\Delta b,\psi_j)$ is the uniform average in
$a$ of $\int_{\log q(a_j)}^{\log q(a)}v_j(q^{-1}(e^\xi))\,d\xi$.
This follows from $b'=l$, $l'=zl$, $d(\log q)/da=d$, and
$f''=l(1+3b+kz)=lh_*$.
Put
\[
 F_j(a)=\mathcal N_j(v_j(a))^2.
\]

We record the bounds furnished by Lemma~\ref{sc:bounds}.  If the endpoint
$a_j$ belongs to $R_1,U_1,V_1,T_1$, respectively, then
\begin{equation}\label{sc:F-simple}
 F_j(a)\le .295,\quad .23,\quad .23,\quad .39
 \qquad\hbox{for all }a.
\end{equation}
For $a\le-1$, the $R_1$ bound improves to $.20$.
For $a_j\in P_1$, the bounds are $.42^2$ when $a\le-1$ and $.61^2$
when $a\ge-1$.  If $a_j\in O$, write $t_j=-a_j$; then
\begin{equation}\label{sc:F-negative}
 \frac{F_j(a)}{\chi(a)^2}\le
 \begin{cases}
 20+1.04(7+t_j^2),&a\le-1,\\
 12+1.63(7+t_j^2),&a\ge-1.
 \end{cases}
\end{equation}
For clarity, the complete finite arithmetic behind these bounds is given
by the following table.  Its three entries bound, in order,
\[
 |W_j|,\qquad
 |A_*b_j-\eta_*/2+\mathsf L_jz+K_{0j}W_j|,\qquad
 |-\eta_*+J_{0j}W_j|.
\]
\begin{equation}\label{sc:substitution-table}
\begin{array}{c|cc}
 &a\le-1&a\ge-1\\\hline
 O&(1.867,.52,1.314)&(2.34,.55,1.155)\\
 P_1&(1.766,.49,1.813)&(2.564,.433,1.155)\\
 R_1&(1.681,.664,2.057)&(2.621,.420,1.155)\\
 U_1&(1.55,.847,2.101)&(2.461,.615,1.155)\\
 V_1&(1.263,.912,1.787)&(1.885,.770,1.205)\\
 T_1&(.913,.869,1.386)&(1.071,.840,1.251)
\end{array}
\end{equation}
To verify it, substitute the intervals in
\eqref{sc:interval-table} and~\eqref{sc:derivative-table}; each expression
is affine in each of its variables separately, so its extrema occur at
corners of the corresponding box.  For the middle expression, first write
it as
\[
 (A_*-2K_{0j})b_j-\eta_*/2+K_{0j}h_*+
 (1-2K_{0j})\mathsf L_jz.
\]
Substituting these three bounds into~\eqref{sc:increment-norm}, together
with the bounds for $\mathsf D_j,\ell_{1j},z,\chi$, proves
\eqref{sc:F-simple}--\eqref{sc:F-negative}.  All these are finite rational
comparisons; the exact verifier accompanying Appendix~\ref{cert:verification}
also checks them.

It remains to compare the integrated derivative cost with the logarithmic
separation surplus.  Let $u_-$ be the mean of
$(\log q(a)-\log q(a_-))/\delta$ under the uniform probability measure
on $[a_-,a_+]$, and put $u_+=1-u_-$.  In the $\xi$ coordinate this measure
has density proportional to $D_0=1/d$.  The excess-variance bound in
Lemma~\ref{tr:q-properties} gives
\[
 0\le \frac{d\log D_0}{d\xi}=\frac{1-qd}{d^2}\le1.
\]
The density $D_0$ is increasing, whereas $D_0e^{-\xi}$ is decreasing.
The covariance of an increasing function with an increasing function is
nonnegative, and with a decreasing function is nonpositive. Applying this
observation to the coordinate $\xi$ compares its mean first with constant
density and then with density proportional to $e^\xi$. The latter comparison
gives
$u_-\le\tfrac12+\tfrac12\coth(\delta/2)-1/\delta$, and hence
\begin{equation}\label{sc:mean-distance}
 \tfrac12\le u_-\le\tfrac12+\delta/12,
 \qquad u_-^2+u_+^2\le\tfrac12(1+\delta^2/36).
\end{equation}
For the upper bound one uses $\coth x-1/x\le x/3$, obtained by comparing
the power series of $x\cosh x$ and $(1+x^2/3)\sinh x$.

Suppose $P_-$ and $P_+$ have the following property: on every partial
$\xi$-interval starting at the indicated endpoint, the mean of
$\mathcal N_j(v_j)$ is at most $\sqrt{P_j}$.  Pointwise bounds for $F_j$
are one way to ensure this property.  Minkowski's inequality, followed by
the uniform average in $a$, gives
\[
 \mathcal N_-(\Delta k,\Delta b,\psi_-)\le\sqrt{P_-}\,\delta u_-,
 \qquad
 \mathcal N_+(\Delta k,\Delta b,\psi_+)\le\sqrt{P_+}\,\delta u_+.
\]
Since $u_-\ge u_+\ge0$,
\[
 P_-u_-^2+P_+u_+^2
 \le\max\left(P_-,\frac{P_-+P_+}{2}\right)(u_-^2+u_+^2).
\]
Averaging~\eqref{sc:one-endpoint} and applying
\eqref{sc:mean-distance} therefore proves the desired estimate whenever
\begin{equation}\label{sc:threshold}
 \max\left(P_-,\frac{P_-+P_+}{2}\right)\le
 \Theta_{m_*}(\delta):=
 \frac{1+m_*}{4}\,
 \frac{1+\delta^2/30+\delta^4/1680}{1+\delta^2/36}.
\end{equation}
Indeed, the resulting coefficient of $Z^2$ is at most
$\delta^2(1+\delta^2/36)\Theta_{m_*}(\delta)/4
=r_{-+}\delta^2$.

We verify~\eqref{sc:threshold} for all endpoint positions.  The function
$\Theta_m$ increases in $\delta\ge0$: differentiation with respect to
$y=\delta^2$ leaves a positive numerator
$1/180+2y/1680+y^2/(1680\cdot36)$.
We will also use the following observation about partial intervals:
if a smaller bound holds on an initial segment of logarithmic length $L$,
then its fraction in a partial interval of length $h\le\delta$ is at least
$\min(1,L/h)\ge L/\delta$, provided $\delta\ge L$.
If $a_-\in R_1$, take $P_-\le.295$ and $P_+\le.39$;
then $(P_-+P_+)/2\le.3425<\Theta_{.38}(0)=.345$.
For $a_-\in U_1$ or $V_1$, use $P_-\le.23$, $P_+\le.39$ and
$\Theta_{.68}(0)=.42$.  For $a_-\in T_1$, both bounds are $.39<.42$.

Now suppose $a_-\in P_1$.  If $a_+\le-1$, both endpoint costs are at
most $.20<\Theta_{.19}(0)$, using the $R_1$ row at $a_+=-1$.
Otherwise, the part of the logarithmic
interval below $-1$ has length at least $1.62$.  Hence every partial
interval from the lower endpoint has mean norm at most
\[
 .61-\frac{(.61-.42)1.62}{\delta},\qquad
 P_-=(.61-.3078/\delta)^2.
\]
The length assertion follows from $q(-1)>.287$, $q(-2)<.056$.
If the upper endpoint lies in $T_1$, then $\delta>4$ because $q(4)>4$.
For $\delta\le4$, therefore, $P_+\le.23$ and
$P_-\le.284143<\Theta_{.19}(0)=.2975$.
For $4\le\delta\le5$, use $P_-<.300787$, $P_+\le.39$, and
\[
 \max(.300787,(.300787+.39)/2)<.345394
 <\Theta_{.19}(4).
\]
For $\delta\ge5$, use $P_-\le.61^2=.3721$ and
\[
 \max(.3721,(.3721+.39)/2)=.38105<\Theta_{.19}(5).
\]
This covers $P_1$.

Finally suppose $a_-\in O$.  For any endpoint $a_j\in O$ and any
argument $a\le-1$ between the endpoints,
\[
 t_j^2\le a^2+2\delta,\qquad |a|\chi(a)\le .09.
\]
The first follows by integrating $d\ge-a$ on the negative half-line;
the second follows from $l\le .075\cdot1.19<.09$ and $d\ge|a|$.
Using $\chi\le.07$ in~\eqref{sc:F-negative} gives
\begin{equation}\label{sc:O-low-cost}
 F_j(a)\le .144+.011\delta.
\end{equation}
If the whole interval lies at or below $-1$, take
$P_-=.144+.011\delta$ and
$P_+\le\max(P_-,.20)$; the latter bound also covers upper endpoints
in $P_1$ and $R_1$.
For $0\le\delta\le4$, both are below $.25$.  For $\delta\ge4$, use
\begin{equation}\label{sc:theta-tail}
 \Theta_0(\delta)\ge .25+.00260\delta^2.
\end{equation}
This follows on clearing denominators and using $\delta^2\ge16$.
The quadratic $.106-.011\delta+.00260\delta^2$ is strictly positive,
so~\eqref{sc:threshold} follows in this case too.

If $a_+>-1$, the first part of the logarithmic interval has length at
least $4$, since $q(-3)<.005$ and $q(-1)>.287$.
For $a\ge-1$, the same integration gives $t_-^2\le1+2\delta$;
using $\chi\le.097$ in~\eqref{sc:F-negative} yields
$F_-(a)\le .236+.031\delta$.
On every partial interval from the lower endpoint, Jensen's inequality
therefore permits the uniform bound
\[
 P_-=.236+.031\delta-\frac{4}{\delta}(.092+.020\delta)
     =.156+.031\delta-.368/\delta.
\]
For partial intervals contained below $-1$, the same bound follows from
\eqref{sc:O-low-cost}, because $\delta\ge4$.
By~\eqref{sc:theta-tail},
\[
 \Theta_0(\delta)-P_-
 \ge .094-.031\delta+.00260\delta^2+.368/\delta>0;
\]
the quadratic has negative discriminant and positive leading coefficient.
If $a_+\notin T_1$, then $P_+\le.23<\Theta_0(\delta)$.
If $a_+\in T_1$, then $\delta>6.5$ from $q(4)>4$ and $q(-3)<.005$,
and
$P_+\le.39<\Theta_0(6.5)<\Theta_0(\delta)$.
This proves~\eqref{sc:threshold} in every case and completes the proof.
\end{proof}
